% Resumo exigido pelo guia PIBIC (máximo de 500 palavras). Resultados marcados como condicionais
% devem ser congelados somente depois da rodada experimental final.
\begin{resumo}
Sistemas de controle que utilizam redes neurais podem produzir ações adequadas em
simulação e, ainda assim, não oferecer uma justificativa verificável para todos os
estados admitidos. Este relatório investiga a aplicabilidade do ESBMC à análise de
propriedades de um controlador neural DDPG para o sistema Cart-Pole, considerando
explicitamente o domínio de entradas, o horizonte de verificação e a representação
numérica utilizada. O ator analisado possui arquitetura 4--24--24--1, com 48
neurônios nas camadas ocultas. Para viabilizar a análise por aritmética inteira,
os pesos e estados são representados em ponto fixo Q8.8, com fator de escala 256,
e os harnesses C são comparados com a implementação correspondente do simulador.

Na rodada canônica de quantização, executada com \texttt{RandomState(42)}, foram
reavaliados 10.000 estados. O erro absoluto médio entre as saídas Float32 e Q8.8
foi de 0,0778 N, o percentil 95 foi de 0,2262 N, o percentil 99 foi de 1,0927 N
e o erro máximo foi de 7,7589 N. O pior estado registrado apresentou inversão
do sinal da ação (+6,4308 N em Float32 e -1,3281 N em Q8.8). A varredura da
aproximação de $\tanh$ encontrou erro absoluto máximo de 0,0457493, equivalente
a 4,5749\% do intervalo unitário. Esses resultados são observações da
representação exportada; a origem no checkpoint PyTorch ainda precisa ser
confirmada diretamente.

Os registros históricos indicam 24 neurônios vivos em cada camada oculta, mas a
auditoria identificou bounds artificiais e independência indevida entre camadas;
por isso, essa análise não é apresentada como prova universal. Na malha fechada,
o harness curto canônico obteve veredicto \texttt{SUCCESSFUL} para o limite de força
$|F|\leq 10$ N com Z3 em 3,5 s,
enquanto Boolector não está disponível no executável Windows. O contraexemplo de
segurança em um passo foi reproduzido como estado fixo, mas esse replay não prova
sozinho a refutação universal. A propriedade direcional permanece inconclusiva
na rodada canônica (timeout com Z3 e solver indisponível com Boolector). A
contribuição é distinguir prova, refutação, reprodução, observação empírica e
resultado inconclusivo, associando cada afirmação ao domínio e ao artefato que a
sustentam.

\textbf{Palavras-chave}: verificação formal; ESBMC; controlador neural; Cart-Pole; quantização Q8.8.
\end{resumo}

\begin{resumo}[Abstract]
\begin{otherlanguage*}{english}
Neural control systems may perform adequately in simulation without providing a
verifiable justification for every admitted state. This report investigates the
applicability of ESBMC to the analysis of a DDPG neural controller for the Cart-Pole
system, explicitly considering the input domain, verification horizon and numerical
representation. The analyzed actor has a 4--24--24--1 architecture, with 48 hidden
neurons. To enable integer-arithmetic analysis, weights and states are represented
in Q8.8 fixed point with scale factor 256, and the C harnesses are compared with
the corresponding simulator implementation.

The canonical quantization run, executed with \texttt{RandomState(42)}, reevaluated
10,000 states. The mean absolute error between Float32 and Q8.8 outputs is
0.0778 N, the 95th percentile is 0.2262 N, the 99th percentile is 1.0927 N,
and the maximum absolute error is 7.7589 N. The worst recorded state reverses
the action sign (+6.4308 N in Float32 and -1.3281 N in Q8.8). A scan of the
$\tanh$ approximation found a maximum absolute error of 0.0457493, or 4.5749\%
of the unit range. These are observations of the exported representations; the
PyTorch checkpoint origin still requires direct confirmation.

Historical records list 24 live neurons in each hidden layer, but the audit found
artificial bounds and improper independence between layers; this analysis is
therefore not presented as a universal proof. In closed-loop analysis, the
canonical short harness proved the force bound $|F|\leq 10$ N with Z3 in 3.5 s,
while Boolector is unavailable in the Windows executable. The one-step safety
counterexample was reproduced as a fixed state, but this replay alone is not a
universal refutation. The directional property remains inconclusive in the
canonical run (Z3 timeout and Boolector unavailable). The contribution is to
distinguish proof, refutation, replay, empirical observation and inconclusive
outcome, linking each claim to its domain and supporting artifact.

\textbf{Keywords}: formal verification; ESBMC; neural controller; Cart-Pole; Q8.8 quantization.
\end{otherlanguage*}
\end{resumo}
